\documentclass[11pt]{article}
\usepackage{geometry}
\geometry{a4paper, margin=1in}
\usepackage{graphicx}
\usepackage{listings}
\usepackage{xcolor}
\usepackage[utf8]{inputenc}
\usepackage[T1]{fontenc}
\usepackage{lmodern}
\usepackage{float}
\usepackage{hyperref}

\lstset{
    basicstyle=\ttfamily\small,
    breaklines=true,
    frame=single,
    language=C++,
    keywordstyle=\color{blue},
    commentstyle=\color{green!50!black},
    stringstyle=\color{red},
    numbers=left,
    numberstyle=\tiny\color{gray},
    showstringspaces=false
}

% Source code lives in ../src, screenshots in ./screenshots
\graphicspath{{screenshots/}}
\newcommand{\src}{../src}

% \screenshot{file}{caption}{label}
% Shows screenshots/<file> if it exists; otherwise a placeholder box so the
% report still compiles before the screenshots are added.
\newcommand{\screenshot}[3]{%
  \begin{figure}[H]
    \centering
    \IfFileExists{screenshots/#1}{%
      \includegraphics[width=\textwidth,height=0.85\textheight,keepaspectratio]{#1}%
    }{%
      \fbox{\parbox[c][4cm][c]{0.9\textwidth}{\centering
        \textbf{Screenshot missing}\\[4pt]
        Save it as \texttt{report/screenshots/\detokenize{#1}}}}%
    }
    \caption{#2}
    \label{#3}
  \end{figure}}


\begin{document}
\sloppy

\title{CS-375 Operating Systems Lab Assignment: Thread-Safe Data Structures}
\author{Santiago Soto} % <-- replace with your name
\date{\today}
\maketitle

\tableofcontents
\newpage

%=====================================================================
\section{Introduction}
This report documents the implementations and analyses for the thread-safe
data structures lab assignment, covering thread-safe queue and stack
implementations, Producer-Consumer and Undo-Redo use cases, and additional
exercises (priority queue, circular buffer, deque and linked list).

All programs were compiled with
\texttt{g++ -o <name> <name>.cpp -pthread -std=c++11} on Ubuntu Linux.
Every data structure follows the same principle: the underlying container is
\emph{private} and every public method acquires the object's
\texttt{std::mutex} before touching it. Because the only way to reach the data
is through a locked method, at most one thread can read or modify the
container at any moment, which removes data races by construction.

All test programs also use a global \texttt{coutMutex} when printing. Without
it, \texttt{std::cout} calls from different threads can interleave
character-by-character and produce garbled lines. This lock is for readable
output only; correctness of the data structures does not depend on it.

%=====================================================================
\section{Exercise 1: Thread-Safe Queue and Stack}
\lstinputlisting[language=C++]{\src/thread_safe_data_structures.cpp}

\subsection*{Explanation}
\textbf{How \texttt{std::lock\_guard} ensures thread safety.}
\texttt{std::lock\_guard<std::mutex> lock(mtx);} locks the mutex in its
constructor and unlocks it in its destructor (RAII). The lock therefore covers
exactly the scope of the method and is released on every exit path, including
the early \texttt{return false} in \texttt{pop()} and any exception thrown by
the container. The mutex cannot be left locked by mistake, which would
deadlock every later caller.

In \texttt{ThreadSafeQueue} and \texttt{ThreadSafeStack}, \texttt{push},
\texttt{pop}, \texttt{empty} and \texttt{size} all lock the same mutex, so they
are mutually exclusive. The key design decision is that \texttt{pop()} does the
\emph{check}, the \emph{read} and the \emph{removal} in one critical section:
\begin{lstlisting}
bool pop(T& value) {
    std::lock_guard<std::mutex> lock(mtx);
    if (queue.empty()) return false;
    value = queue.front();
    queue.pop();
    return true;
}
\end{lstlisting}
If a caller instead wrote \texttt{if (!q.empty()) \{ v = q.front(); q.pop(); \}}
using separately locked calls, another thread could empty the queue between
the check and the read (a \emph{check-then-act} race). Combining the steps
makes the operation atomic.

\textbf{Producer-Consumer.} Three producers each push 5 messages into a shared
\texttt{ThreadSafeQueue<std::string>}, and two consumers pop them. The
counter \texttt{messages\_produced} is a \texttt{std::atomic<int>}, so its
increments from several producers are never lost.

\textbf{Consumer termination condition.} A consumer loops forever and stops
only when \texttt{pop()} fails \emph{and} all 15 messages had been produced. A
failed \texttt{pop()} alone is not enough, because the queue may just be
temporarily empty while producers are still running. Since each producer
increments the counter \emph{after} the push, a counter value of 15 means all
15 messages are already in the queue. The order of the checks matters: the
consumer reads the counter \emph{before} calling \texttt{pop()}. If it checked
afterwards, this interleaving would be possible: the pop fails, a producer pushes
the last message and sets the counter to 15, then the consumer sees 15 and quits
with one message still in the queue. Reading first guarantees that a failed pop
really means the queue has been drained. (The original starter
loop \texttt{while (consumed < 15)} counted per consumer, so with two consumers
neither one would ever reach 15 by itself. The \texttt{break} was the only real
exit, and I made that explicit.) The summary line confirms
\texttt{produced 15, consumed 15, left in queue 0}.

\textbf{Undo-Redo.} \texttt{editText} pushes the current text onto the undo
stack and clears the redo stack, because a new edit invalidates the redo
history. \texttt{undo} pops the previous state and pushes the current state onto
the redo stack. \texttt{redo} does the reverse. The sequence produces:
Hello $\rightarrow$ Hello World $\rightarrow$ Hello Universe $\rightarrow$
(undo) Hello World $\rightarrow$ (undo) Hello $\rightarrow$ (redo) Hello World
$\rightarrow$ Hello Galaxy. A final \texttt{redo} prints ``Nothing to redo'',
which shows that the edit cleared the redo history.

\subsection*{Analysis}
Output order varies between runs because the OS scheduler decides which thread
runs next. That variation is expected. What never varies is that each message
is consumed \emph{exactly once}, and none are lost or duplicated.

\textbf{If synchronization were removed:} \texttt{std::queue} is not
thread-safe. Two producers calling \texttt{push} simultaneously could both
write to the same slot or corrupt the internal \texttt{std::deque} block
pointers, which can lose messages or crash the program. Two consumers could both
read the same \texttt{front()} before either calls \texttt{pop()}, so one
message is processed twice and the next is popped without ever being read. If
\texttt{messages\_produced} were a plain \texttt{int}, the read-modify-write in
\texttt{++} could lose increments. The counter might then never reach 15, and
the consumers would spin forever.

\textbf{Deadlock considerations:} each method holds exactly one lock and never
calls another locking method while holding it, so no lock-ordering cycle (and
no deadlock) is possible. In the Undo-Redo system, each stack operation is
individually atomic. The \emph{combination} of \texttt{undoStack.pop} +
\texttt{redoStack.push} + updating \texttt{currentText} is not atomic, however.
That is acceptable here because the demo is single-threaded. With several
editing threads, one extra mutex around each whole edit/undo/redo would be
needed to keep the three pieces of state consistent.

\textbf{Observed output (Figure~\ref{fig:ex1}).} The three producers' messages
interleave in a different order in each round (for example Producer 0, 2, 1 in
round 1), which shows the scheduler running the threads truly concurrently.
Consumer 0 processed 8 messages and Consumer 1 processed 7, so the work was
shared between them, and the summary line \texttt{produced 15, consumed 15,
left in queue 0} confirms that every message was consumed exactly once. Both
consumers printed their ``finished'' line, which shows the termination
condition works. In Problem 2, the stack states follow the expected sequence,
and the final \texttt{Redo - Nothing to redo} shows that the ``Hello Galaxy''
edit cleared the redo history (final undo stack size 3, redo stack size 0).

\subsection*{Screenshot}
\screenshot{Lab7exercise1_screenshot.png}{Compilation and execution of thread\_safe\_data\_structures.cpp}{fig:ex1}

%=====================================================================
\section{Exercise 3: Thread-Safe Priority Queue}
\lstinputlisting[language=C++]{\src/priority_queue.cpp}

\subsection*{Explanation}
\textbf{Maintaining order.} \texttt{std::priority\_queue<int>} is a binary
max-heap stored in a \texttt{std::vector}. \texttt{push} appends the element
and sifts it up, and \texttt{pop} moves the last element to the root and sifts
it down. Both take $O(\log n)$ and preserve the heap invariant (each parent
$\ge$ its children), so \texttt{top()} is always the largest element.

\textbf{Thread safety.} During a sift the heap is temporarily \emph{invalid}.
Elements are halfway through being swapped. If another thread read or modified
the heap at that moment, it could observe or create a broken heap, after which
\texttt{top()} would no longer be the maximum. Holding the mutex for the whole
\texttt{push}/\texttt{pop} makes each heap update atomic, so other threads only
ever see a valid heap. As in Exercise 1, \texttt{pop} combines \texttt{top()}
and \texttt{pop()} under one lock.

I replaced \texttt{rand()} with a per-thread \texttt{std::mt19937} generator
because \texttt{rand()} uses hidden global state and is not guaranteed to be
thread-safe. The popper now loops until it has removed all 20 elements instead
of making a fixed 20 attempts.

\subsection*{Analysis}
Four pusher threads push 5 random priorities each (every 100\,ms). One popper
removes an element every 150\,ms. While pushes are still arriving, the popper
returns the largest element \emph{currently} in the queue. A larger value
pushed later can therefore appear after a smaller one was popped earlier, so
the popped sequence is not globally sorted. Once all pushers finish (about
0.5\,s), the remaining pops come out in strictly descending order, which
demonstrates that the heap order is preserved under concurrency. The final line
confirms all 20 elements were popped and the queue is empty.

Because every operation takes the same lock, pushers and the popper are fully
serialized on the queue. With heavy contention this lock becomes a bottleneck,
but each critical section is only $O(\log n)$, so contention here is minimal.

\textbf{Observed output (Figure~\ref{fig:ex3}).} The first pop returned 88,
but larger values (98, 93, 99) were pushed later. While pushers are still
running, each pop returns the maximum of what is \emph{currently} in the queue,
not of everything that will ever be pushed. Once all 20 values were in the
queue, the popper drained it in strictly descending order (99, 89, 88, 87,
\ldots, 20, 20, 10), and the duplicate value 20 was handled correctly. The
\texttt{remaining} counter drops by exactly one per pop down to 0, so no
element was lost or duplicated.

\subsection*{Screenshot}
\screenshot{Lab7exercise3.png}{Compilation and execution of priority\_queue.cpp}{fig:ex3}

%=====================================================================
\section{Exercise 4: Thread-Safe Circular Buffer}
\lstinputlisting[language=C++]{\src/circular_buffer.cpp}

\subsection*{Explanation}
The buffer is a fixed array of \texttt{BUFFER\_SIZE = 5} integers with a write
index \texttt{in}, a read index \texttt{out} and a \texttt{count}. Both indices
advance with \texttt{(i + 1) \% BUFFER\_SIZE}, so they wrap around and reuse
slots without any allocation.

\textbf{Condition variables.} Besides mutual exclusion, a bounded buffer needs
\emph{waiting}: a producer must wait while the buffer is full and a consumer
must wait while it is empty. Two \texttt{std::condition\_variable}s provide
this:
\begin{itemize}
  \item \texttt{not\_full.wait(lock, [this]\{ return count < BUFFER\_SIZE; \})}
        atomically releases the mutex and puts the producer to sleep. It
        re-acquires the lock when woken and returns only once the predicate is
        true.
  \item After inserting, the producer calls \texttt{not\_empty.notify\_one()}
        to wake one sleeping consumer. Consumers do the symmetric thing with
        \texttt{not\_empty} / \texttt{not\_full}.
\end{itemize}
\texttt{std::unique\_lock} is required instead of \texttt{lock\_guard} because
\texttt{wait()} must unlock and relock the mutex. The predicate form of
\texttt{wait} re-checks the condition in a loop, which handles
\emph{spurious wake-ups} and the case where another thread consumed the slot
first. The lock is released before \texttt{notify\_one()} so the woken thread
does not immediately block on the mutex again.

\textbf{Change from the starter code.} In the starter code, \texttt{pop()}
returned \texttt{false} on an empty buffer instead of waiting. The consumer's
fixed 10-iteration loop could then finish having consumed fewer than 10
items. Items would be silently left in the buffer, and if more than 5 were
left the producer would block on \texttt{not\_full} forever (a hang). I made
\texttt{pop()} wait on \texttt{not\_empty}, as the requirement to use condition
variables implies, and kept the non-blocking behaviour as a separate
\texttt{try\_pop()}. The test uses 2 producers and 2 consumers (10 items per
producer, 10 per consumer), as required.

\subsection*{Analysis}
Producers run every 50\,ms and consumers every 150\,ms, so the buffer fills to
5/5 (printed as \texttt{FULL}). Producers then block until a consumer frees a
slot. This is \emph{back-pressure}: the fast side is throttled to the speed of
the slow side, and memory use stays bounded. All 20 items are consumed and the
buffer ends empty.

\textbf{Comparison with \texttt{ThreadSafeQueue} (Exercise 1):}
\begin{center}
\begin{tabular}{|p{3.3cm}|p{5.5cm}|p{5.5cm}|}
\hline
 & \textbf{ThreadSafeQueue} & \textbf{ThreadSafeCircularBuffer} \\ \hline
Capacity & Unbounded (grows with \texttt{std::queue}) & Fixed, 5 slots \\ \hline
Memory & Dynamic allocation on growth & Pre-allocated array, no allocation \\ \hline
Synchronization & \texttt{mutex} + \texttt{lock\_guard} only & \texttt{mutex} + two condition variables \\ \hline
Empty pop & Returns \texttt{false} immediately; caller must poll (sleep and retry) & Blocks until data arrives, with no busy-waiting \\ \hline
Full push & Never full & Blocks until space is free \\ \hline
CPU use while idle & Polling wastes wake-ups & Sleeping threads use no CPU \\ \hline
Risk & Unbounded memory if producers outpace consumers & Deadlock if consumers stop while producers wait \\ \hline
\end{tabular}
\end{center}
Without the condition variables, a producer would have to spin while the
buffer is full, wasting CPU. Without the mutex, two producers could write the
same \texttt{buffer[in]} slot and both increment \texttt{count}, which loses
one item and corrupts the count.

\textbf{Observed output (Figure~\ref{fig:ex4}).} From the moment
\texttt{104} was produced, the buffer repeatedly reaches \texttt{5/5 FULL}.
After that, every new item appears only right after a consumer frees a slot,
which is the producers blocking on \texttt{not\_full} as intended. All 20 values
(100--109 and 200--209) are consumed exactly once, and the run ends with
\texttt{Buffer empty: yes}.

Some printed occupancies look odd, for example \texttt{Producer 2 produced:
200 [buffer 0/5]} or \texttt{Consumer 2 consumed: 203 [buffer 5/5]}. This is
not a bug in the buffer. The program calls \texttt{size()} \emph{after}
\texttt{push()}/\texttt{pop()} has released the lock, so other threads can run in
between. In the first case a consumer had already removed item 200; in the
second a producer had already refilled the slot. Each method is atomic on its
own, but a \emph{sequence} of calls is not, so any value read after the lock is
released is only a snapshot. This is exactly the check-then-act issue
discussed in Exercise 1.

\subsection*{Screenshot}
\screenshot{Lab7exercise4_screenshot.png}{Compilation and execution of circular\_buffer.cpp}{fig:ex4}

%=====================================================================
\section{Exercise 5: Thread-Safe Deque}
\lstinputlisting[language=C++]{\src/thread_safe_deque.cpp}

\subsection*{Explanation}
\textbf{Challenges of synchronizing both ends.} It is tempting to give the
front and the back their own locks so a thread at one end does not block a
thread at the other. The two ends are not independent, however:
\begin{itemize}
  \item When the deque holds \textbf{one} element, \texttt{front()} and
        \texttt{back()} are the \emph{same} element. A concurrent
        \texttt{pop\_front} and \texttt{pop\_back} could both see it and both
        remove it. The result is one value returned twice and an underflow of
        the container.
  \item When it is \textbf{empty}, a \texttt{push\_front} and a
        \texttt{pop\_back} race on the same element.
  \item \texttt{std::deque} keeps shared bookkeeping (block map, size). Either
        end can reallocate the block map, which invalidates the other end's
        internal pointers.
\end{itemize}
For these reasons the implementation uses \textbf{one mutex for the whole
deque}. Each operation, including the \texttt{empty()} check plus the
read and removal in the pops, is a single atomic critical section. Fine-grained
per-end locking is possible only with a custom structure (for example a doubly
linked list with sentinel nodes), and it still needs special handling when the
size is 0 or 1, which adds complexity and room for deadlock if two locks are
ever taken in different orders.

\subsection*{Analysis}
Thread 1 pushes to the front, thread 2 pushes to the back, and two threads
pop from the front and back. Elements from thread 1 accumulate at the front
(most recent first) and elements from thread 2 at the back, so
\emph{Popped front} mostly returns 1xx values and \emph{Popped back} mostly
returns 2xx values. When the deque gets short, either popper can take an
element pushed from the other end. The final summary checks that
\texttt{pushed = popped + remaining}, so no element was lost or duplicated.
Some pop attempts find the deque empty and simply return \texttt{false}.

Without the mutex, the size-0/1 races above would lead to double pops,
lost elements or a crash. The cost of the single lock is that operations on
opposite ends cannot run in parallel. For short critical sections like these,
that cost is negligible compared to the safety gained.

\textbf{Observed output (Figure~\ref{fig:ex5}).} Front pops returned 1xx
values and back pops returned 2xx values, as expected, since thread 1 only
pushes at the front and thread 2 only at the back. One line appears out of
order: \texttt{Popped back: 200} is printed before \texttt{Thread 2 pushed back:
200}. The push really did happen first (otherwise the pop could not have
returned 200). The pushing thread released the deque lock and was then
preempted before it acquired \texttt{coutMutex}, so the popper's message was
printed first. Print order is therefore not the same as operation order. The
summary \texttt{pushed 10, popped 10, remaining 0 -> consistent} confirms that
no element was lost or duplicated.

\subsection*{Screenshot}
\screenshot{Lab7exercise5_screenshot.png}{Compilation and execution of thread\_safe\_deque.cpp}{fig:ex5}

%=====================================================================
\section{Exercise 6: Thread-Safe Linked List}
\lstinputlisting[language=C++]{\src/thread_safe_linked_list.cpp}

\subsection*{Explanation}
\textbf{Maintaining list integrity.} A singly linked list is defined entirely
by its \texttt{head} pointer and the \texttt{next} pointers. Each operation
modifies them in several steps, and the list must never be observed or modified
between those steps:
\begin{itemize}
  \item \textbf{Lost insert:} \texttt{push\_front} reads \texttt{head} into
        \texttt{newNode->next} and then writes \texttt{head = newNode}. If two
        threads interleave, both new nodes point to the same old head, and the
        second write overwrites the first. One node becomes unreachable
        (lost data and a memory leak).
  \item \textbf{Double delete:} if two threads run \texttt{pop\_front}
        simultaneously, both may read the same \texttt{head}, return the same
        value and \texttt{delete} the same node twice. That is undefined
        behaviour, typically heap corruption or a crash.
  \item \textbf{Use after free:} \texttt{size()} walks the list. If a node is
        deleted while another thread is standing on it, the walker dereferences
        freed memory.
\end{itemize}
The mutex makes every read-modify-write of \texttt{head} atomic, and
\texttt{size()} holds the lock for the whole traversal. Two small optimizations
shorten the critical sections without breaking safety. \texttt{new Node} is
allocated \emph{before} taking the lock because it touches no shared state, and
\texttt{delete temp} happens \emph{after} releasing the lock because the node
is already unlinked and no other thread can reach it. I also deleted the copy
constructor and copy assignment, because a copied list would share and
double-delete the same nodes. The destructor pops all remaining nodes so no
memory leaks.

\subsection*{Analysis}
Two pushers insert 5 nodes each at the head while a popper removes from the
head. Since this is a LIFO structure, the popper usually gets the most recently
pushed node. After pushing stops, nodes come out in reverse insertion order.
The list size printed after each pop decreases consistently.

\textbf{Performance considerations.} This is \emph{coarse-grained} locking:
one mutex serializes all operations. It is simple and obviously correct.
\texttt{push\_front}/\texttt{pop\_front} are $O(1)$, but \texttt{size()} is
$O(n)$ and blocks every other thread for the whole traversal. Possible
improvements:
\begin{itemize}
  \item keep a \texttt{count} member updated under the lock, so \texttt{size()}
        is $O(1)$;
  \item \emph{fine-grained (hand-over-hand) locking} with one mutex per node,
        which lets operations on different parts of the list proceed in
        parallel. It is more complex, and locks must always be acquired in list
        order to avoid deadlock;
  \item a \emph{lock-free} stack using \texttt{std::atomic<Node*>} and
        compare-and-swap on \texttt{head}. This has no blocking, but it must
        deal with the ABA problem and safe memory reclamation.
\end{itemize}
Since all operations here act only on the head, fine-grained locking would
gain nothing. A single mutex is the right choice for this workload.

\textbf{Observed output (Figure~\ref{fig:ex6}).} Because
\texttt{push\_front}/\texttt{pop\_front} form a LIFO stack, each pop returns
the most recently inserted node (for example 100 right after it was pushed). After
both pushers finished, the remaining nodes came out in reverse insertion order
(204, 104, 203, 202, 102, 201, 200). The reported list size decreases by one with
each pop down to 0, and all 10 nodes were removed exactly once. The list
was never corrupted (no lost node, no double free, no crash).

\subsection*{Screenshot}
\screenshot{Lab7exercise6_screenshot.png}{Compilation and execution of thread\_safe\_linked\_list.cpp}{fig:ex6}

%=====================================================================
\section{Conclusion}
All five data structures achieve thread safety the same way: private state, one
mutex per object, and RAII locks (\texttt{std::lock\_guard} /
\texttt{std::unique\_lock}) that cover each operation entirely, including
compound \emph{check-then-act} steps such as ``if not empty, read and remove''.
The circular buffer adds condition variables so threads can \emph{wait} for a
state change instead of polling. Removing the locks would produce lost updates,
duplicated or lost elements, corrupted containers, double frees and crashes.
Keeping each critical section short and never holding two locks at once
prevents deadlock.

\end{document}
